\documentclass{article}

\usepackage{amsmath}
\usepackage{amssymb}
\usepackage{geometry}
\geometry{margin=1in}

\title{Research on dbf design docs}
\author{Author}
\date{\today}

\begin{document}

\maketitle

\section{2025 UNSW}

\begin{enumerate}
	\item 3.1.3. They named their subsystem requirements for easy access later on, they gave them labels. 
	\item 3.1.3. They have their subsystem requirements and missio overview requiremets together in one table whereas we have it in two tables. 
	\item 3.3. They determined a maximum from previous score, at least the performance data
	\item 3.3. Assume maximum score, and then deduct form there for sesitivity? Explain why sensitivity matters. 
	\item 3.3. Sensitivity score analysis did not yield any conclusions about how we should try to score. 
	\item 3.4. Their Aircraft configuration decisions were very clear on how they came up to the decisio, they listed steps to how they decided it.
	\item 3.4. They put absolutely everything into that design stage, the weighing and stuff.
	\item 4.2. They had a whole model for the flight and described how the calculations validated each part of the flight
	\item 4.3.1. They had a lot of formulas describing the thrust to weight ratios
	\item 4.3.2. They used the specs of the location of flight and explained how everything was used
	\item 4.3.2. They compared to previous planes to check reasonability of their plane
	\item 4.3.2. They varied one component and then picked the maximum of that one, and showed the plot for it, checked feasibility
	\item 4.3.4. For taper, they made a graph varying tapers and then showed the resulting plots
\end{enumerate}

Software/equations to look into: Gudmundsson, Matlab

They had a sensitivity study for every part of the design basically. 

\end{document}